\section{System and Threat Model}
\label{sec:model}

\subsection{Roles}
Figure~\ref{fig:arch} shows the roles. A provider-run \emph{orchestrator tier} (orchestrator, regional routers,
operational database, attestation verifier) is trusted for availability and accounting but never sees tenant secrets.
Each tenant runs its own \emph{key broker} (KBS). \emph{Workers} are volunteer machines booted into the attested
\voss{} appliance, which runs tenant functions in Firecracker microVMs. One worker per cluster additionally plays the
\emph{scheduler} role; every attested worker carries the scheduler code in the same image, so promoting a successor costs
no boot. The orchestrator partitions workers into small geographic clusters (tens of workers), which bounds the planner's
problem size and intra-cluster latency.

\begin{figure}[t]
\centering
\resizebox{\columnwidth}{!}{%
\begin{tikzpicture}[
  node distance=5mm and 7mm,
  box/.style={draw, rounded corners=2pt, minimum height=7mm, align=center, font=\scriptsize, inner sep=2.5pt},
  trusted/.style={box, fill=blue!7},
  attested/.style={box, fill=green!9},
  tenant/.style={box, fill=orange!12},
  arr/.style={-{Latex[length=1.6mm]}, thick},
  dsh/.style={-{Latex[length=1.6mm]}, dashed}]
  \node[box] (client) {Client};
  \node[trusted, right=of client] (router) {Router\\[-1pt]\tiny(per region)};
  \node[attested, right=of router] (sched) {Scheduler\\[-1pt]\tiny(role in one worker)};
  \node[attested, right=of sched] (worker) {Workers\\[-1pt]\tiny(appliance + microVMs)};
  \node[trusted, above=7mm of router] (orch) {Orchestrator + DB};
  \node[trusted, above=7mm of sched] (ver) {Verifier};
  \node[tenant, above=7mm of worker] (kbs) {Tenant KBS};
  \draw[arr] (client) -- node[above,font=\tiny]{invoke} (router);
  \draw[arr] (router) -- node[above,font=\tiny]{enqueue} (sched);
  \draw[arr] (sched) -- node[above,font=\tiny]{dispatch} (worker);
  \draw[dsh] (orch) -- node[left,font=\tiny]{clusters,} node[right,font=\tiny, xshift=-9mm, yshift=-1.5mm]{succession} (sched);
  \draw[dsh] (worker) -- node[right,font=\tiny]{quote} (ver);
  \draw[dsh] (worker) -- node[right,font=\tiny]{key} (kbs);
  \draw[dsh] (router) -- node[left,font=\tiny]{receipts} (orch);
  \node[draw, dotted, fit=(sched)(worker), inner sep=3pt, label={[font=\tiny]below:one cluster (tens of workers)}] {};
\end{tikzpicture}}
\caption{Roles. Blue: provider-run orchestrator tier (trusted for availability and accounting only). Green:
volunteer machines running the attested appliance (their \emph{hosts} may still omit or delay). Orange: tenant-operated
key broker. Solid arrows: the invocation path; dashed: control and trust.}
\label{fig:arch}
\end{figure}

\subsection{Donation sessions and attestation}
A worker alternates between \textsc{off} and \emph{sessions}. A session starts when the owner reboots into the
appliance (scheduled, idle-triggered, or on a dedicated spare box), becomes schedulable when attestation completes, and
ends abruptly (owner reboot, power loss) or gracefully (a drain request with a deadline). Session lengths are heavy-tailed
and time-of-day dependent, with regional shocks. The appliance image is a signed unified kernel image with a dm-verity
root, no owner login, and a data partition sealed to the expected PCR values; the disk caches survive a session as
ciphertext.

Attestation follows the Keylime pattern \cite{schear2016keylime}. The agent presents the TPM's endorsement-key
certificate, then a quote over PCRs and the IMA log bound to a nonce and a freshly generated session public key; the
verifier replays the event log against a policy and returns a session certificate that expires unless re-quoted. All
channels use mutual TLS under that certificate, so no untrusted OS lies on any path. A tenant's KBS releases that tenant's
key, wrapped to the session key, only to such a host. \emph{Trust-warm} means that the key is present in appliance memory
this session; it is lost with the session. The scheduler and orchestrator never see tenant keys.

\subsection{Threat model}
\label{sec:threat}
\emph{Volunteer owners} control power, network, firmware settings and their own OS, may try to boot modified software,
run the appliance in a VM, replay, delay or drop traffic, create identities, and choose when to defect. Measured boot, a
vendor-certified TPM endorsement key, an owner-locked image and sealed storage defeat the first group; what remains is
\emph{omission and timing}: power off, drop, delay, throttle, and lying about anything outside the attested stack
(including fake drain notices). \emph{Network attackers} see only ciphertext. \emph{Other tenants} are contained by
microVM isolation. Out of scope: lab-grade physical attacks (for example fTPM voltage glitching \cite{jacob2023faultpm}),
kernel and VMM zero-days in the appliance, side channels and traffic analysis, and failure of or attacks on the
orchestrator tier. Identity is the hash of the TPM endorsement key, so identities cost machines \cite{douceur2002sybil}.

\textbf{Consequence.} To the rest of the system a malicious worker or scheduler host is an omission/timing-faulty node.
The same mechanisms therefore cover churn, crashes and malice (Sec.~\ref{sec:design}).

\subsection{Substrate profiles}
\label{sec:substrates}
A \emph{substrate profile} is data that tells one generic lifecycle engine how a worker establishes trust, how long it
lasts, what the warm-state ladder looks like, and what assurance a host of that class carries. We model five.
\texttt{appliance\_tpm} (primary) pays a reboot, measured boot and a TPM quote once per session. \texttt{cvm\_snp}
and \texttt{coco\_pool} model server confidential VMs and confidential containers with per-wCVM or per-shell establishment and
trustlet-style warm levels, with parameters from the literature (ranges swept in experiments). \texttt{drtm\_hv} models
a measured hypervisor launched under the owner's OS (no reboot, owner interference). \texttt{none} is the unattested reference.

For appliance and unattested hosts the ladder has five levels per (worker, function): L0 nothing cached; L1 code bundle and
root file system on disk; L2 snapshot on disk; L3 warm idle VM in RAM; L4 busy. Trust is an orthogonal flag $K[t,j]$.

\subsection{Exposure}
Let $a_j\in[0,1]$ be the assurance of host $j$ (substrate base, freshness of the last quote, firmware policy), and
$\mathrm{KH}(t)$ the set of hosts holding tenant $t$'s key or decrypted artifacts. The \emph{exposure} of $t$ is
$E_t=\sum_{j\in \mathrm{KH}(t)}(1-a_j)$, weighted by the tenant's sensitivity class $\sigma_t$ in the planner. It is a
transparent proxy for the attack surface, not a probability of compromise, and we therefore report latency, exposure and
robustness side by side rather than as one weighted number.
